Many local community service nonprofits report difficulty employing enough staff to meet rising demand for services, with rising burnout rates among workers. For those employed in this sector in food pantries, homeless shelters, child social work, the question is if the rising trend in artificial intelligence (AI) could potentially offload some of the worker burden. I examine this question for the three largest occupations in community food, housing, and relief services in the United States, combining three published task-level measures of theoretical exposure to AI, time spent per task, and observed AI use. Two of the occupations studied showed that the amount of time that large language models (LLMs) could directly speed up amounts to about a half an hour per day. A greater share of the workday (between 28-55 percent of the day) could be alleviated if additional software were built on top of the LLM for the purpose of supporting things like developing care plans and implementing training programs. These findings suggest that general-purpose AI tools are likely to return little frontline staff time, and that larger gains depend on purpose-built tools that extend LLM capability.
